\bookchapter{บทที่ 17: การเชื่อมต่อมือถือกับอุปกรณ์อื่น}

\chapterintro{มือถือในมือคุณแท้จริงแล้วคือรีโมตสุดฉลาดผสมศูนย์กลางข้อมูล มันไม่ใช่แค่โทรได้ แต่ยังเชื่อมคอมพิวเตอร์ ทีวี นาฬิกา ลำโพง และรถยนต์เข้าด้วยกัน ทำให้อุปกรณ์เหล่านี้ไม่ใช่เกาะเดี่ยวอีกต่อไป บทนี้เราจะคุยกันว่าจะทำให้พวกมัน “รวมทีมจัดการงาน” กันได้อย่างไร}

\booksection{หนึ่ง การทำงานร่วมกับคอมพิวเตอร์: มือถือกับคอมไม่ใช่สองโลกอีกต่อไป}

พูดตามตรง สิ่งที่ผมรำคาญที่สุดเมื่อก่อนคือถ่ายรูปด้วยมือถือแล้วอยากย้ายขึ้นคอมไปแต่งต่อ ต้องหาสายเคเบิล รอถ่ายโอน หลังจากนั้นค้นพบว่า จริง ๆ ไม่ต้องใช้สายเคเบิลเลย การเชื่อมต่อที่ดีเร็วกว่าสายเคเบิลสิบเท่า

\begin{itemize}
  \item \textbf{ระบบนิเวศ Apple (AirDrop + คลิปบอร์ดร่วม)}: ถ้าคุณใช้ iPhone คู่กับ Mac การส่งแอร์ดรอป (AirDrop) คือฟังก์ชันแรกที่คุณควรเรียนรู้ บนมือถือกดแชร์ AirDrop เลือก Mac ของคุณ รูปหายไปในพริบตา ไม่ใช้เน็ต ไม่ใช้ Wi-Fi แต่ใช้บลูทูธกับ Wi-Fi เชื่อมตรงกัน ยังมีฟังก์ชันที่คนมองข้ามง่ายชื่อ "คลิปบอร์ดร่วม (Universal Clipboard)" — คัดลอกข้อความบนมือถือ แล้ววางบน Mac ได้เลย เหมือนเวทมนตร์ชะมัด
  \item \textbf{Android + Windows (เชื่อมต่อถึง Windows)}: Microsoft มีแอปชื่อ "เชื่อมต่อถึง Windows" มือถือ Samsung ติดตั้งมาในเครื่อง ยี่ห้ออื่นไปโหลดจากร้านแอปหนึ่งตัวก็ใช้ได้ จับคู่แล้วคุณสามารถดูการแจ้งเตือนมือถือบนคอมได้ตรง ๆ ตอบข้อความสั้น ส่องรูป แม้กระทั่งรับ-โทรออก ผมใช้มาสองปี ตอนประชุมไม่ต้องคอยก้มดูมือถือ เห็นการแจ้งเตือนเด้งบนจอคอมก็พอ
  \item \textbf{การทำงานร่วมกันหลายจอของแบรนด์}: "การทำงานร่วมกันหลายจอ" ของ Huawei และ HyperConnect ของ Xiaomi เถือกกว่านั้น — ฉายหน้าจอมือถือขึ้นคอมได้ตรง ๆ ใช้เมาส์ลากไฟล์จากมือถือไปคอมได้ ยังดูวิดีโอสั้น ตอบ WeChat บนคอมได้ คัดลอกวางข้ามอุปกรณ์ ลากไฟล์ก็ทำได้หมด ประสบการณ์ลื่นมาก
\end{itemize}

\booksection{สอง แคสต์ไปที่ทีวี: ทั้งครอบครัวดูด้วยกัน}

ตอนกลับบ้านเทศกาลปีใหม่ น้องสาวลูกพี่ลูกน้องอยากฉายวิดีโอทริปรับปริญญาในมือถือลงทีวี งงอยู่ครึ่งวัน สุดท้ายค้นพบว่าจริง ๆ ง่ายนิดเดียว:

\begin{itemize}
  \item \textbf{iPhone ทีวี}: ปัดลงจากมุมขวาบนเปิดศูนย์ควบคุม กดไอคอนที่เป็นกรอบสองอันซ้อนกัน (การฉายหน้าจอ) หาชื่อทีวีของคุณ กดปุ๊บเชื่อมต่อปั๊บ เงื่อนไขคือทีวีกับมือถือต้องต่อ Wi-Fi เดียวกัน และทีวีต้องรองรับ AirPlay (ทีวีอัจฉริยะส่วนใหญ่รองรับ)
  \item \textbf{Android ทีวี}: แต่ละบ้านเรียกไม่เหมือนกัน Huawei เรียก "การฉายจอไร้สาย" Xiaomi เรียก "การแคสต์" OPPO/vivo เรียก "การร่วมมือหลายจอ" ปกติหาได้ในแผงสวิตช์ลัด ถ้าไม่เจอไปค้นหาคำว่า "แคสต์" ในการตั้งค่า
  \item \textbf{แอปวิดีโอแคสต์ตรง ๆ}: ถ้าคุณแค่อยากดู BiliBili, YouTube บนทีวี จริง ๆ ง่ายกว่านั้น — ที่หน้าเล่นหาไอคอนเล็ก ๆ ที่เป็นสี่เหลี่ยมมีสัญลักษณ์ Wi-Fi กดหนึ่งครั้ง เลือกทีวี วิดีโอก็ไปแล้ว ตอนนั้นมือถือทำอย่างอื่นต่อได้ วิดีโอเล่นต่อบนทีวี ไม่รบกวนกัน
  \item \textbf{แคสต์ไม่สำเร็จทำอย่างไรดี}: 90\% ของกรณีคือไม่ได้ต่อ Wi-Fi เดียวกัน อีก 10\% คือฟังก์ชันแคสต์ของทีวียังไม่เปิด ไปที่การตั้งค่าทีวีค้นหา "การร่วมมือหลายจอ" หรือ "AirPlay" แล้วเปิดเสีย
\end{itemize}

\booksection{สาม กับสมาร์ทวอทช์/สายรัดข้อมือ: ผู้จัดการสุขภาพอยู่บนข้อมือคุณ}

ผมซื้อสายรัดข้อมือให้แม่ เธอตอนแรกบ่นว่า “ฉันไม่ได้ออกกำลังกาย ซื้อทำนะ” ผลคือใช้ไปสองสัปดาห์ เธอพบว่าตัวเองนอนได้แค่วันละห้าชั่วโมงกว่า ก็เริ่มปรับเวลานอนเอง นี่แหละคุณค่าที่ใช้ได้จริงที่สุดของนาฬิกา/สายรัด — เปลี่ยนข้อมูลสุขภาพที่คุณมองไม่เห็นให้เป็นตัวเลขที่มองเห็นได้

\begin{itemize}
  \item \textbf{ตัวกรองการแจ้งเตือน}: เวลามือถืออยู่ในกระเป๋าหรือกระเป๋ากางเกงสั่นไม่หยุด ยกข้อมือดูหนึ่งตา — คือ WeChat ของเจ้านายหรือโปรโมชันของแอปช้อปปิ้ง ตัดสินใจได้ในหนึ่งวินาที ข้อความที่ไม่คุ้มหยิบมือถือก็ปล่อยมันไป
  \item \textbf{ซิงก์ข้อมูลสุขภาพ}: นาฬิกาวัดอัตราการเต้นหัวใจ ออกซิเจนในเลือด การนอน และจำนวนก้าวตลอด 24 ชั่วโมง ข้อมูลพวกนี้จะรวบรวมเข้าแอปสุขภาพบนมือถืออัตโนมัติ ผ่านไปนาน ๆ คุณจะเห็นเทรนด์ของตัวเอง — เดือนนี้เดินน้อยไปหรือเปล่า สัปดาห์นี้คุณภาพการนอนดีขึ้นไหม
  \item \textbf{ซิงก์บัตรโดยสาร/บัตรเข้าออก}: บัตรโดยสารและบัตรเข้าออกที่ผูกไว้บนมือถือ แบรนด์ใหญ่ ๆ ส่วนใหญ่จะซิงก์ขึ้นนาฬิกาให้อัตโนมัติ ตอนนั่งรถไฟฟ้าไม่ต้องหยิบมือถือ เอาข้อมือแตะประตูก็ผ่าน มีความรู้สึกแบบ "อนาคตมาแล้ว" จริง ๆ
  \item \textbf{ของวิเศษหามือถือ}: บนนาฬิกาปกติมีฟังก์ชัน "ค้นหามือถือ" กดหนึ่งครั้งมือถือก็จะส่งเสียง ผมใช้ฟังก์ชันนี้ถี่เกินคาด — มือถือวางปุ๊บลืมว่าวางไว้ตรงไหนเนี่ย
\end{itemize}

\booksection{สี่ กับสมาร์ทโฮม: มือถือคือ “สวิตช์ใหญ่” ของบ้าน}

สมาร์ทโฮมฟังดูลี้ลับ แต่เอาเข้าจริงมีแค่สามเรื่อง: ควบคุมจากระยะไกลด้วยมือถือ ตั้งกฎอัตโนมัติ และสั่งงานด้วยเสียง

เล่าการใช้งานจริงไม่กี่อย่างของผมเอง คุณจะรู้ทันทีว่าของพวกนี้ใช้ได้จริงแค่ไหน:

\begin{itemize}
  \item \textbf{สมาร์ทโฮมฟังดูลี้ลับ แต่เอาเข้าจริงมีแค่สามเรื่อง: ควบคุมจากระยะไกลด้วยมือถือ ตั้งกฎอัตโนมัติ และสั่งงานด้วยเสียง}เล่าการใช้งานจริงไม่กี่อย่างของผมเอง คุณจะรู้ทันทีว่าของพวกนี้ใช้ได้จริงแค่ไหน:
  \item \textbf{เครื่องปรับอากาศ}: วันร้อนกลับบ้าน เปิดแอร์ไว้ล่วงหน้าจากมือถือก่อนถึงบ้าน พอเข้าประตูก็เย็นฉ่ำแล้ว จ่ายเงินร้อยกว่าซื้อปลั๊กอัจฉริยะเสียบเข้าไปก็ทำได้ ไม่ต้องเปลี่ยนแอร์
  \item \textbf{ไฟ}: ก่อนนอนนอนบนเตียง ขี้เกียจลุกไปปิดไฟ กดที่มือถือหนึ่งครั้งหรือสั่งหนึ่งเสียง "เฮ้ Siri / ผู้ช่วยเสียง ปิดไฟ"
  \item \textbf{หุ่นยนต์ดูดฝุ่น}: ก่อนออกจากบ้านกด "เริ่มกวาด" ที่มือถือหนึ่งครั้ง กลับมาพื้นก็สะอาดแล้ว ยังตั้งให้กวาดเวลาทุกวันได้ด้วย
\end{itemize}

กล้องวงจรปิด

\begin{itemize}
  \item : บ้านมีคุณปู่คุณย่า เด็กเล็ก หรือสัตว์เลี้ยง เปิดดูสถานการณ์ในบ้านจากมือถือได้ตลอด
  \item เลือกระบบนิเวศไหนดี? สรุปสั้น ๆ หนึ่งประโยค:
  \item : ใช้ iPhone ดีที่สุดเลือกอุปกรณ์ที่รองรับ Apple HomeKit (แพงนิดหน่อยแต่ปกป้องความเป็นส่วนตัวดี)
\end{itemize}

\booksection{หูฟังบลูทูธและลำโพง: อิสรภาพไร้สาย}

: ใช้มือถือ Huawei ดีที่สุด HarmonyOS Connect ประสบการณ์กับอุปกรณ์ Huawei ดีที่สุด

\begin{itemize}
  \item \textbf{ห้า หูฟังบลูทูธกับลำโพง: อิสรภาพไร้สาย}จากหูฟังมีสายเปลี่ยนมาเป็นหูฟังบลูทูธ สำหรับผมคือการอัปเกรดประสบการณ์มือถือที่เห็นความต่างชัดที่สุดเรื่องหนึ่ง ไม่ต้องควักก้อนเส้นพันกันออกจากกระเป๋าอีก ก็ไม่ต้องกลัวเส้นเกี่ยวที่มือจับประตู
  \item \textbf{จับคู่ง่ายมาก}: หูฟังเปิดเครื่องครั้งแรกจะเข้าโหมดจับคู่อัตโนมัติ (ไฟกระพริบ) เปิดบลูทูธบนมือถือค้นหาได้เลย กดหนึ่งครั้งก็เชื่อมต่อ หลังจากนั้นหยิบออกมาเมื่อไรก็เชื่อมอัตโนมัติ ไม่ต้องทำอะไรเลย
  \item \textbf{แบรนด์เดียวกันมีโบนัส}: AirPods จับคู่ iPhone จะป๊อปอัปแสดงแบตอัตโนมัติ และสลับใช้ระหว่างอุปกรณ์ Apple หลายตัวได้ไร้รอยต่อ Huawei FreeBuds จับคู่มือถือ Huawei หูฟัง Xiaomi จับคู่มือถือ Xiaomi ก็มีป๊อปอัปและการเพิ่มประสิทธิภาพคล้ายกัน
  \item \textbf{ลำโพงบลูทูธ}: งานรวมญาติ แคมป์กลางแจ้ง ลำโพงบลูทูธหนึ่งตัวช่วยเพิ่มบรรยากาศได้เท่าตัว เชื่อมต่อแล้ว มือถือคือแผง DJ ของคุณ
\end{itemize}

\begin{tcbinfo}{เทคโนโลยี UWB (Ultra-Wideband) — "การรับรู้พื้นที่" ที่แม่นยำกว่า Bluetooth}
นอกจาก Bluetooth แล้ว ในมือถือรุ่นใหม่ยังมีชิป UWB (Ultra-Wideband) ซ่อนอยู่อีก 1 ตัว ความสามารถของมันคือ\textbf{การวัดระยะทางและการระบุตำแหน่งในพื้นที่อย่างแม่นยำ}—สามารถคำนวณได้ว่าอุปกรณ์อีกเครื่องอยู่ห่างจากคุณแค่ไหนและอยู่ทิศทางใด โดยมีความแม่นยำระดับเซนติเมตร Bluetooth บอกได้แค่ว่า “อยู่ใกล้ๆ” แต่ UWB สามารถบอกได้ว่า “อยู่ด้านหน้าซ้ายของคุณ 1.2 เมตร” การใช้งานทั่วไป:

\begin{itemize}
  \item \textbf{ค้นหา AirTag อย่างแม่นยำ}: เมื่อ iPhone ค้นหา AirTag จะแสดงลูกศรนำทางว่า "อยู่ห่างไป 7 เมตร เดินไปทางซ้าย" ซึ่งใช้ UWB ในขณะที่ Bluetooth ทั่วไปบอกได้แค่ว่า "อยู่ในระยะหรือไม่"
  \item \textbf{กุญแจรถดิจิทัล}: iPhone และมือถือ Android เรือธงที่รองรับ UWB สามารถทำหน้าที่เป็นกุญแจรถ เข้าใกล้รถจะปลดล็อกอัตโนมัติ เดินออกห่างจะล็อกรถอัตโนมัติ และยังสามารถแยกแยะตำแหน่งคนขับและผู้โดยสารได้ เพื่อป้องกันการแฮ็กแบบ Relay
  \item \textbf{นำทางในอาคาร}: ในห้างสรรพสินค้าหรือสนามบิน สามารถนำทางระดับ "ทางเดิน" ได้ผ่าน UWB Beacon ซึ่งแม่นยำกว่า GPS มาก
  \item \textbf{รุ่นที่รองรับ}: iPhone 11 ขึ้นไปมาพร้อมชิป UWB (U1/U2) ทั้งหมด ส่วนฝั่ง Android บางรุ่นเรือธงของ Xiaomi, Samsung, vivo และ OPPO ก็เริ่มติดตั้งแล้ว
\end{itemize}

\end{tcbinfo}

\begin{tcbinfo}{LE Audio และ Auracast — มาตรฐานใหม่ของเสียง Bluetooth}
เสียง Bluetooth แบบดั้งเดิม (A2DP) ที่ใช้กันมายี่สิบปีกำลังถูกแทนที่โดย\textbf{LE Audio}(เสียงพลังงานต่ำ) ซึ่งเป็นมาตรฐานเสียงรุ่นใหม่ตั้งแต่ Bluetooth 5.2 เป็นต้นไป:

\begin{itemize}
  \item \textbf{ประหยัดไฟกว่า คุณภาพเสียงดีกว่า}:ใช้การเข้ารหัส LC3 ใหม่ล่าสุด ซึ่งให้คุณภาพเสียงที่ดีกว่าและใช้พลังงานน้อยกว่ามาตรฐาน SBC เดิม ทำให้แบตเตอรี่หูฟังใช้งานได้นานขึ้นอย่างเห็นได้ชัด
  \item \textbf{เสียงสองช่องทาง}:ส่งสัญญาณหูฟังซ้ายและขวาแยกจากกัน สลับใช้โหมดหูฟังข้างเดียวโดยไม่หลุดการเชื่อมต่อ
  \item \textbf{Auracast กระจายเสียง}:ฟีเจอร์เด่นของ LE Audio โทรศัพท์/ทีวีเครื่องเดียวสามารถกระจายเสียงไปพร้อมกัน\textbf{ไม่จำกัด}หูฟัง — ในห้องรอเครื่องบิน พิพิธภัณฑ์ ยิม หรือห้องเรียน สามารถใช้แหล่งส่งสัญญาณเดียว ให้หลายคนใช้หูฟังของตัวเองฟังเสียงจากภาพเดียวกันได้
  \item \textbf{สถานะปัจจุบัน}:หูฟังและโทรศัพท์เรือธงที่วางตลาดใหม่ในช่วง 2024-2026 จะรองรับ LE Audio เป็นส่วนใหญ่ อุปกรณ์รุ่นเก่าต้องอัปเกรดฮาร์ดแวร์ ไม่สามารถปลดล็อกผ่านซอฟต์แวร์ได้ เวลาซื้อหูฟังใหม่สามารถสังเกตว่ามีโลโก้รองรับ LE Audio/Auracast หรือไม่
\end{itemize}

\end{tcbinfo}

\booksection{หก กับระบบในรถยนต์: คู่หูความปลอดภัยตอนขับรถ}

ขับรถแล้วดูมือถือผิดกฎหมาย และเสี่ยงต่อชีวิต แต่บางทีก็จำเป็นต้องนำทาง รับสาย เปิดเพลง ตอนนั้น CarPlay (ฝั่ง Apple) กับ Android Auto (ฝั่ง Android) ก็เข้ามามีประโยชน์

\begin{itemize}
  \item \textbf{วิธีเชื่อมต่อ}: ถ้ารถคุณรองรับ (รถใหม่ช่วงหลัง ๆ ส่วนใหญ่รองรับ) ใช้สายเคเบิลเสียบมือถือเข้ากับพอร์ต USB บนรถ หน้าจอในรถจะสลับเป็นหน้าจอ CarPlay หรือ Android Auto อัตโนมัติ รถรุ่นใหม่บางรุ่นรองรับเชื่อมต่อไร้สาย ยิ่งสบาย
  \item \textbf{ใช้ทำอะไรได้}: นำทาง (ใช้ดีกว่านำทางในรถเยอะ) โทรออก (รับสายที่ปุ่มพวงมาลัย) ฟังเพลง (แอปเพลงบนมือถือส่งตรงเข้ารถ) ส่งข้อความ WeChat ด้วยเสียง (ให้ Siri หรือผู้ช่วยเสียงช่วยอ่านข้อความและตอบ)
  \item \textbf{รถเก่าก็ใช้ได้}: ถ้ารถคุณเก่าเกินไปไม่รองรับ จ่ายหลักร้อยซื้อจอ CarPlay แบบติดตั้งเสริมมาแปะหน้าปัด ประสบการณ์ต่างกันไม่มาก
\end{itemize}

\booksection{เจ็ด ความปลอดภัยและความเป็นส่วนตัวในการเชื่อมต่อ}

อุปกรณ์เชื่อมต่อยิ่งมาก พื้นที่เสี่ยงรั่วไหลของความเป็นส่วนตัวก็ยิ่งกว้างขึ้น มีไม่กี่จุดที่ต้องระวัง:

\begin{itemize}
  \item : ก่อนฉายจอในที่สาธารณะ ดูแถบการแจ้งเตือนมือถือก่อนว่ามีข้อความส่วนตัวไหม อย่าให้ทั้งบ้านเห็นของที่ไม่ควรเห็น
  \item : ก่อนส่งต่ออุปกรณ์ให้คนอื่นหรือขายทิ้ง ยกเลิกการผูกก่อน — ไปที่การตั้งค่ามือถือหาอุปกรณ์นั้น เลือก "นำออก" หรือ "ยกเลิกการจับคู่"
  \item : กล้องวงจรปิดสมาร์ทโฮม ซื้อแบบมีแบรนด์ มีการเข้ารหัส อย่าซื้อของแปะแบรนด์ราคาถูก ๆ ข่าวกล้องถูกแฮกเกอร์เอาไปถ่ายทอดสดก็ได้ยินกันมาเยอะแล้ว
\end{itemize}

\begin{tcbpractice}{มาลองดู: ฉายหน้าจอมือถือขึ้นทีวี แล้วตั้งค่าสถานการณ์อัจฉริยะหนึ่งอย่าง}
\begin{enumerate}
  \item : ยืนยันว่ามือถือกับทีวีต่อ Wi-Fi เดียวกัน
  \item : บนทีวีหาฟังก์ชัน "การฉายหน้าจอ" หรือ "AirPlay" แล้วเปิด
  \item : บนมือถือ: iPhone เปิดจากศูนย์ควบคุมกด "การฉายหน้าจอ" ฝั่ง Android ปัดลงค้นหา "การฉายจอไร้สาย"
  \item : เลือกรูปที่คุณชอบที่สุดหนึ่งรูปฉายขึ้นทีวี ให้คนในบ้านดูดู
  \item : ถ้ามีอุปกรณ์สมาร์ทโฮม ลองสร้าง "โหมดออกจากบ้าน" ในแอปบนมือถือ: กดปุ๊บปิดไฟ ปิดแอร์ เริ่มงานหุ่นยนต์ดูดฝุ่น
\end{enumerate}

\end{tcbpractice}

\begin{tcbinfo}{เคส 1: การปรับแต่งแบตเตอรี่ให้ใช้งานได้นานขึ้น}
\begin{itemize}
  \item \textbf{สถานการณ์}:ผู้ใช้ Xiaomi 14 แบตเตอรี่ไม่ค่อยทน ใช้ไปครึ่งวันก็หมด
  \item \textbf{ขั้นตอน}:(1) ตั้งค่า แบตเตอรี่ ดูอันดับแอปที่กินไฟ (2) ระบุแอปที่ทำงานเบื้องหลังผิดปกติ (3) ปิด 5G สลับไปใช้ 4G (4) จำกัดการทำงานเบื้องหลัง (5) เปิดโหมดประหยัดแบตเตอรี่ (6) ถอดเคสระบายความร้อนตอนชาร์จ
  \item \textbf{ผลลัพธ์ที่คาดหวัง}:แบตเตอรี่ใช้งานได้นานขึ้น 30-50\% และสามารถอ่านข้อมูลสุขภาพแบตเตอรี่ได้
\end{itemize}

\end{tcbinfo}
